\chapter{Расшифровка критериев оценки (баллы 0--3)}
\label{app:criteria-rubric}
\index{критерии оценки!расшифровка}

Актуальная таблица уровней и путей к файлам поддерживается в репозитории: \path{docs/criteria\_rubric.md}. Номинальная сумма \textbf{raw} по автоматической части --- до \textbf{75}.

Каждый критерий \textbf{C01--C25} оценивается по уровням \textbf{0 / 1 / 2 / 3} (в т.\,ч. \textbf{C14} --- до 3 баллов за согласованное размещение \texttt{numpy} в SBOM\_TCB/SBOM\_OTHER). Экспертные \textbf{C20--C22} в скрипте дают \textbf{0}; жюри выставляет баллы вручную.

\section*{Автоматический скрипт}
Реализация: \path{scripts/evaluate\_contest\_score.py}. Особенности:
\begin{itemize}
  \item \textbf{C02:} все тесты решения должны находиться в \texttt{src\_solution/tests/**}.
  \item \textbf{C03:} маркер \texttt{security} в \texttt{pytest.ini} и использование в тестах.
  \item \textbf{C12:} подсчёт файлов в \texttt{tests/}, где AST содержит \texttt{import src\_solution}, а не подстрока в комментарии.
  \item \textbf{C13:} раздел тестов безопасности в отчёте решения в каталоге \texttt{src\_solution/docs/} с явными путями \texttt{src\_solution/...}.
  \item \textbf{C15:} тесты с импортами и \texttt{event\_log}, и \texttt{src\_solution}.
\end{itemize}

\section*{C01--C19 (кратко)}
Полный текст см.\ \path{docs/criteria\_rubric.md}. Уровни согласованы с версией скрипта в репозитории (релиз документации \textbf{1.7.7}).

\section*{C20: Стоимость сертификации (рейтинг)}
Выставляется \textbf{жюри} после анализа стоимости сертификации у всех участников (значения из логов/отчётов Регулятора).
\begin{itemize}
  \item \textbf{3} --- наименьшая стоимость среди конкурсантов.
  \item \textbf{2} --- второе место по стоимости.
  \item \textbf{1} --- третье место.
  \item \textbf{0} --- все остальные участники по этому критерию.
\end{itemize}
В автоматическом скрипте критерий даёт \textbf{0}; итог вносится в отчёт вручную.

\section*{C21: Экспертно --- соответствие политик архитектуре АБУ}
\begin{itemize}
  \item \textbf{0--3} --- жюри оценивает согласованность политик безопасности, реализованных в решении, с архитектурой АБУ и моделью угроз (в скрипте \textbf{0}).
\end{itemize}

\section*{C22: Экспертно --- полнота отчёта и воспроизводимость}
\begin{itemize}
  \item \textbf{0--3} --- жюри оценивает полноту отчёта, воспроизводимость команд, ясность аргументации (в скрипте \textbf{0}).
\end{itemize}

\section*{C23--C25: Дополнительные критерии по ДВБ и монитору}
\begin{itemize}
  \item \textbf{C23} --- размер доверенных доменов (максимальный LOC одного домена ДВБ).
  \item \textbf{C24} --- число разрешающих IPC-политик для доверенного домена.
  \item \textbf{C25} --- наличие \texttt{security\_monitor}, security-тестов и порогового покрытия monitor-кода.
\end{itemize}
